\begin{defi}

\label{def:TA}
A \emph{table algebra} (TA) $(A, B)$ is a finite dimensional
algebra $A$
over the complex numbers $\C$, and a distinguished basis $B$ that
contains $1_A$, such that the following properties hold:
\begin{enumerate}[label=(1\alph*)]
\item\label{defi1.1_1a} %[(1a)]
The structure constants for $B$ are all nonnegative real numbers; that is, for all
$b,c \in B$,
\[
bc = \sum_{d \in B} \lambda_{bcd}d, \quad \text{for some }
\lambda_{bcd} \in \R_{\geq 0}.
\]
\item\label{defi1.1_1b} %[(1b)]
There is an algebra anti-automorphism (denoted by $^*$) of $A$ such
that $(a^*)^* = a$ for all $a \in A$; and $B^* = B$.
\item\label{defi1.1_1c} %[(1c)]
For all $b,c \in B$, 
$\lambda_{bc1} = 0$ if and only if $c \neq b^*$. 
\end{enumerate}
\end{defi}

%\

It follows as a consequence of the definition that $\lambda_{bb^*1} =
\lambda_{b^*b1} > 0$ for all $b \in B$. Frobenius-Perron eigenvalue
theory yields that for each table algebra there exists a unique algebra
homomorphism $\delta: A \rightarrow {\C}$, called the \emph{degree map},
such that $\delta(b) = \delta(b^*) > 0$ for all $b \in B$. The table
algebra $(A, B)$ is called \emph{standard} if for all $b \in B$,
$\delta(b) = \lambda_{bb^*1}$. Any table algebra can be \emph{rescaled}
(replace each $b \in B$ by $\beta_bb$, for suitable $\beta_b > 0$) to
one that is standard.

%\

Let $(A, B)$ be a standard table algebra (STA). For any subsets $S, T$
of $B$, the \emph{set product} $ST \coloneqq \cup_{s\in S, t\in
T}\Supp _B(st)$, $S^* \coloneqq \{s^* | s \in S\}$, and $S^+ \coloneqq \sum_{s
\in S}s$. Note that set product is associative. If $T = \{t\}$, a
singleton set, then $S\{t\}$ (resp. $\{t\}S, S\{t\}S$) is denoted $St$
(resp. $tS, StS$). The \emph{order} of a subset $S$ is $o(S) \coloneqq \delta
(S^+)$.

A nonempty subset $C \subseteq B$ is called \emph{closed} if $C^*C
\subseteq C$. In this case, $({\C}C, C)$ is again a table algebra, and
the set of \emph{left cosets} $Cb$ for $b \in B$ partition $B$, as do
the \emph{right cosets} $bC$, and the $C$-$C$ \emph{double cosets}
$CbC$. A \emph{quotient element}, for any $b \in B$, is $b//C \coloneqq 
o(C)^{-1}(CbC)^+$. Let $B//C \coloneqq \{b//C | b \in B\}$, and
$A//C \coloneqq {\C}(B//C)$. Then $(A//C, B//C)$ is a
STA, called the \emph{quotient algebra}, or \emph{double coset algebra}
of $(A, B)$ by the closed subset $C$. Its degree map is
$\delta\downarrow _{A//C}$, and its anti-automorphism is
${^*}\downarrow_{A//C}$. Furthermore, $o(B//C) =
o(B)/o(C)$. The closed subsets $D$ with $C \subseteq D \subseteq B$ are
in bijection with the closed subsets of $B//C$ via $D
\leftrightarrow D//C$ (see Proposition~\ref{prop:subset_corresp} below). The closed subset $C$ is called \emph{normal} (resp.
\emph{strongly normal}) in $B$ if $bC = Cb$ (resp. $bCb^* = C$) for all
$b \in B$. Strongly normal closed subsets are normal, but the converse
is not always true. 

An element $x \in B$ is called \emph{thin} (or \emph{linear}, or
\emph{grouplike}) if $xx^* = 1$. This is equivalent to $\delta(x) = 1$; 
and if $x$ is thin, then $xb \in B$ and $bx \in B$ for all $b \in B$.
So $\delta(xb) = \delta(bx) = \delta(b)$ for all $b \in B$. Now
$O_{\vartheta}(B)\coloneqq \{x \in B | xx^* = 1\} = \{x \in B |
\delta(x) = 1\}$ is a group, called the \emph{thin radical} of $B$; and
$B$ is called thin if $B = O_{\vartheta}(B)$. Obviously,
$O_{\vartheta}(B)$ is the unique maximal closed subset $D$ of $B$
such that $D$ is thin. For $C$ a closed subset of $B$, the quotient
$B//C$ is thin if and only if $C$ is strongly normal in $B$. Since the
intersection of strongly normal closed subsets is again strongly normal,
there is a unique minimal closed subset of $B$, denoted
$O^{\vartheta}(B)$, such that $B//O^{\vartheta}(B)$ is thin. Now
$O^{\vartheta}(B)$ is called the \emph{thin residue} of $B$, and it
equals the closed subset of $B$ generated by $\Supp _B(bb^*)$ for
all $b \in B$. (See~\cite[Theorem~2.3.1]{Z1} or~\cite [Theorem~3.2.1]{Z2}, 
where the algebraic proof for association schemes holds
verbatim for table algebras.) 

%\

Throughout, $Z(A)$ will denote the center of the algebra $A$, and $Z(B)$
will mean $B \cap Z(A)$.

\begin{defi}
\label{def:chain}
A \emph{chain $\mathcal{C}$ of length} $n$ is a collection $\{C_i\}_{i=0}^n$
of closed subsets of $B$ such that either $C_0 \subset C_1 \subset
\cdots \subset C_{n-1} \subset C_n$ or $C_n \subset C_{n-1} \subset
\cdots \subset C_1 \subset C_0$. 
 It is called a \emph{bi-anchored chain}
(\emph{B-chain}) if $C_0 = \{1\}$ and $C_n = B$ (or $C_n = \{1\}$ and
$C_0 = B$). It is a \emph{thin chain} (\emph{T-chain}) if $C_{i+1}//C_i$
is thin for all $0 \leq i \leq n-1$ (or $C_{i-1}//C_i$ is thin for all
$1 \leq i \leq n$). It is called a \emph{thin-central chain}
(\emph{TC-chain}) if $C_{i+1}//C_i$ is thin and $C_{i+1}//C_i \subseteq
Z(B//C_i)$ for all $0 \leq i \leq n-1$ (or $C_{i-1}//C_i$ is thin and
$C_{i-1}//C_i \subseteq Z(B//C_i)$ for all $1 \leq i \leq n$).
\end{defi}

\begin{defi}
\label{def:res_thin_nilp}
A STA $(A, B)$ is \emph{residually thin} if there exists a bi-anchored
thin chain (BT-chain). It is \emph{nilpotent} if there exists a
bi-anchored thin-central chain (BTC-chain).
\end{defi}

\begin{rema}
\label{rem:nilp_group}
It is immediate from the definition that for any finite group $G$, the
group algebra $({\C}G, G)$ is nilpotent as a STA if and only if $G$ is
nilpotent in the usual group-theoretic sense.
\end{rema}

\begin{prop}
\label{prop:res_thin_order}
Let $(A, B)$ be a residually thin STA, with BT-chain $B = C_0 \supset
C_1 \supset \cdots \supset C_{n-1} \supset C_n = \{1\}$. Then $o(C_i)$
is an integer, and $o(C_{i+1}) \vert o(C_i)$ for all $0 \leq i \leq
n-1$.
\end{prop}

\begin{proof}
Since each $C_j//C_{j+1}$ is a group, $o(C_j//C_{j+1}) =
o(C_j)/o(C_{j+1})$ is an integer for $0 \leq j \leq n-1$. Since $o(C_i)
= o(C_{i+1})o(C_i//C_{i+1}) = \prod_{j \geq i}o(C_j//C_{j+1})$, the
result follows.
\end{proof}

\begin{defi}
\label{def:card_number}
Let $(A, B)$ be a STA with a B-chain $\mathcal{C}$: $B = C_0 \supset C_1
\supset \cdots C_{n-1} \supset C_n = \{1\}$. For any $1 \leq i \leq n$
and any $b \in
B\backslash C_i$, define the positive integer $m(\mathcal{C}, i,
b)$ for $i < n$ by 
\[
m(\mathcal{C}, i, b) \coloneqq \prod_{j=i}^{n-1}\card (C_j//C_{j+1}
\cdot b//C_{j+1} \cdot C_j//C_{j+1}),
\]
where $C_j//C_{j+1} \cdot b//C_{j+1} \cdot C_j//C_{j+1}$ is the
$C_j//C_{j+1}$ - $C_j//C_{j+1}$ double coset of $b//C_{j+1}$ in
$B//C_{j+1}$; and $m(\mathcal{C}, n, b) \coloneqq \card (C_nbC_n) = \card\{b\} = 1$.
\end{defi}
